% ============================================================
% 分片：s03d-twisted  对应原文 PDF 第 35--43 页（印刷页码 35--43），末段延至第 44 页前 9 行
% 负责章节：3.4 The Twisted Models In Detail
% 译者：t07-s34   审校状态：初稿
% 审校：t05-s32 已完成逐式核对（修正 3 处）
% 边界说明（Lead 裁定）：p35 开头两段为 §3.3 的收尾，按 SPEC 页码表归本分片，首句自 p34 末行补全；
%   3.4 末段延至 p44 前 9 行（至 “induce a reparametrization of the variable t.”），§4 标题起归 s04-geometry。
% ============================================================

% ==== 原文 PDF 第 35 页（印刷页 35）====
\label{op:35}
% [跨页接续: 原文第34页末至第35页首] 以下首句自 p34 末行 “The A model has a greatly enriched variant in which one ignores ...” 起，补全为完整句。
A 模型有一个内容大为丰富的变体：其中忽略 $Q_-$，而把 $\bar Q_+$ 取作 BRST 算符；于是物理态或顶点算符就是 $\bar Q_+$ 的上同调类。（半扭模型的物理态空间是无穷维的；椭圆亏格（elliptic genus）(3.29) 本质上是 $\bar Q_+$ 算符的指标。）$A$ 模型的这个变体称为半扭模型（half-twisted model）。半扭模型不是拓扑场论；若把它表述在黎曼面 $\Sigma$ 上，其关联函数是共形不变的（即使未扭模型并不共形不变），并且随 $\Sigma$ 的复结构全纯地变化。依上述论证，半扭模型的关联函数在 $e$ 改变下不变，并同时随 $W$ 与 $t$ 全纯地变化。
% [Tutor Note: 椭圆亏格作为 $\bar Q_+$ 的指标：由于取指标后只保留 $\bar Q_+$-闭态、$\bar q$ 依赖抵消，结果全纯依赖于 $\tau$（$q=e^{2\pi i\tau}$）；$q$ 的指数为 $L_0-c/24$，而 $R$ 荷通过 $y=e^{2\pi iz}$ 的幂（$J_0$ 本征值）进入系数，故椭圆亏格的 $q$-展开是按左行能量 $L_0$（以及 $y$ 幂所记的左行 $R$ 荷 $J_0$）分次的计数——这正是它能作为半扭模型物理态数目的原因。]

由于 Calabi--Yau 与 Landau--Ginzburg 之间的插值是通过改变 $t$ 得到的，可以得出若干结论：$B$ 模型的拓扑可观测量在 Calabi--Yau 与 Landau--Ginzburg 下相同；而对 $A$ 模型与半扭模型，Landau--Ginzburg 是 Calabi--Yau 的解析延拓。

\subsection{扭模型详述}

下面更仔细地考察 $B$ 模型与 $A$ 模型。（以下可以把半扭模型与 $A$ 模型放在一起考虑。）基本工具是 [24,§5] 中讨论的一类不动点定理。回忆 $A$ 模型或 $B$ 模型都具有一个 $(0|2)$ 维的超群 $F$ 作为对称性，其费米型生成元记为 $Q_\sigma$。所有顶点算符都可选为 $F$ 不变的代表元，因此 $F$ 可视为路径积分的对称性。
% [Tutor Note: $(0|2)$ 维指 0 个玻色型、2 个费米型生成元，故指标 $\sigma$ 取两个值。这里的 $Q_\sigma$ 已是拓扑扭变（topological twist）后的 BRST 荷：扭变把原本按世界面 $SO(2)$ 张量变换的超荷重新定义为标量，使能动张量成为 $Q$-恰当项，$Q^2=0$ 且 $T_{\alpha\beta}=\{Q,\Lambda_{\alpha\beta}\}$，关联函数因而与度规无关；$A$、$B$ 两个模型只是选取不同的标量超荷组合，可观测量即相应 $Q$ 的上同调类（$A$ 模型与 $B$ 模型各自的上同调在 §3.3 讨论）。]

拓扑关联函数的路径积分可以局域化到 $F$ 的不动点空间 $Z$ 上，$Z$ 即场空间中满足 $\{Q_\sigma,\Lambda\}=0$（对每个场 $\Lambda$ 与每个 $\sigma$）的点集。若 $Q_\sigma$ 沿 $Z$ 有单零点，则对 $Z$ 的法向模式作路径积分时，由于玻色子与费米子之间的抵消，只给出一个简单的因子 $\pm1$。否则，路径积分归结为场空间中某个有限维方向空间上的积分，在这些方向上 $Q_\sigma$ 的零点阶数高于一阶。
% [Tutor Note: 上述局域化正是不动点定理的路径积分表述：把 $Q_\sigma$ 视为 BRST 算符、在 $\{Q_\sigma,\Lambda\}=0$ 的解集 $Z$ 上作半经典展开，法向方向的高斯积分给出 $\mathrm{Ber}(Q_\sigma^2|_{N})^{-1/2}$，玻色/费米行列式之比为 $\pm1$（超行列式的符号由零模数决定）。这与 Duistermaat--Heckman 型局域化公式同构，也是 [24,§5] 中“拓扑 $\sigma$ 模型的关联函数可化为靶空间上的有限维积分”的来源。]

% ==== 原文 PDF 第 36 页（印刷页 36）====
\label{op:36}
\subsubsection*{$B$ 模型}

在 $B$ 模型中，考察 $\{\bar Q_\pm,\bar\lambda_\pm\}=0$，可得 $D=v_{12}=0$。\footnote{回忆 Wick 转动 $y^0\to-iy^2$，故 $v_{01}\to-iv_{12}$。} 设 $\Phi_\alpha$ 是手征超场，其分量为 $\phi_\alpha,\psi_\alpha,F_\alpha$，荷为 $Q_\alpha$。令 $\{\bar Q_\pm,\bar\psi_{\alpha\pm}\}=0$，可得 $F=0$。由 $\{\bar Q_\pm,\psi_{\alpha\pm}\}=0$ 可知 $D\phi_\alpha/Dy^i=0$，且对一切 $\alpha$ 有 $Q_\alpha\phi_\alpha\sigma=0$。无论 $r\gg0$ 还是 $r\ll0$，方程 $D=F_\alpha=0$ 都要求这些量 $Q_\alpha\phi_\alpha$ 不全为零，于是最后一个方程给出 $\sigma=0$。
% [原文疑误: 由 $\{\bar Q_\pm,\bar\psi_{\alpha\pm}\}=0$ 通常应得 $\bar F=0$（$\bar F$ 为反手征多重态的最高分量辅助场），此处图像中印作无横线的 $F=0$，疑为原稿漏排横线；物理上终态两者皆为零，故照录原文，不作改动。]
% [Tutor Note: 推导补足——这里用到 D 项方程 $D=-e^2(\sum_\alpha Q_\alpha\bar\phi_\alpha\phi_\alpha-r)$（见第 37 页）。由 $D=0$ 得 $\sum_\alpha Q_\alpha|\phi_\alpha|^2=r$；当 $r\gg0$ 时右端为正、当 $r\ll0$ 时右端为负，两种情形下都不可能全部 $Q_\alpha\phi_\alpha=0$，故“不全为零”成立。再与 $Q_\alpha\phi_\alpha\sigma=0$（逐个 $\alpha$）联立即得 $\sigma=0$。]

总之，对 $B$ 模型，$F$ 不动点就是从世界面 $\Sigma$ 到经典真空空间的常值映射。当 $r\gg0$ 时，经典真空空间是低能 $\sigma$ 模型的靶空间 $X$。此时生成元 $Q_\sigma$ 沿其零点空间 $Z\cong X$ 有单零点，因此路径积分归结为 $X$ 上的积分。该积分可按 [24,§4] 的方式分析；关联函数可用 $X$ 上微分形式的周期来求值。当 $r\ll0$ 时，经典真空空间是一个点，故 $Q_\sigma$ 只有一个零点，但该零点是退化的；路径积分归结为 Vafa [31] 分析过的有限维积分。
% [Tutor Note: 物理图像——$B$ 模型的可观测量正是未扭模型的 (cc) 手征环：$Q_\sigma$-闭算符对应 $H^{p,q}(X)$ 中的元素，环的乘法（OPE 的奇点部分）对应 Hodge 结构的上积，而“周期”给出把闭链映为复数的线性泛函。手征环的乘积结构（即三维 Calabi--Yau 上的 Yukawa 耦合）正是由这些周期确定的，这也是 $B$ 模型完全由复结构模控制的原因。]

\subsubsection*{$A$ 模型}

对 $A$ 模型，方程 $\{Q_\sigma,\lambda_\pm\}=0$ 与 $\{Q_\sigma,\bar\lambda_\pm\}=0$ 给出
\begin{equation}
\begin{aligned}
D+v_{12} &= 0\\
\mathrm{d}\sigma &= 0.
\end{aligned}
\tag{3.33}
\end{equation}

方程 $\{Q_\sigma,\psi_{\alpha\pm}\}=0$ 与 $\{Q_\sigma,\bar\psi_{\alpha\pm}\}=0$ 类似地给出
\begin{equation}
\begin{aligned}
\left(\frac{D}{Dy^1}+i\frac{D}{Dy^2}\right)\phi_\alpha &= 0\\
F_\alpha &= 0\\
Q_\alpha\phi_\alpha\sigma &= 0.
\end{aligned}
\tag{3.34}
\end{equation}

式 (3.34) 中最后一个方程与式 (3.33) 中最后一个方程合起来表明 $\sigma=0$，因为对所关心的模型，其余方程不允许 $Q_\alpha\phi_\alpha$ 对每个 $\alpha$ 都恒等于零。
% [跨页接续: 原文第36页末至第37页首；p36 末句 “...do not permit $Q_\alpha\phi_\alpha$ to” 在 p37 首行续为 “vanish identically for each $\alpha$.”]
% ==== 原文 PDF 第 37 页（印刷页 37）====
\label{op:37}
于是幸存的方程是
\begin{equation}
\begin{aligned}
\left(\frac{D}{Dy^1}+i\frac{D}{Dy^2}\right)\phi_\alpha &= 0\\
F_\alpha &= 0\\
D+v_{12} &= 0.
\end{aligned}
\tag{3.35}
\end{equation}

由恒等式
\begin{equation}
\int \mathrm{d}^2y\,D_i\bar\phi_\alpha D^i\phi_\alpha
= \int \mathrm{d}^2y\,(D_1-iD_2)\bar\phi_\alpha\,(D_1+iD_2)\phi_\alpha
- \int \mathrm{d}^2y\,Q_\alpha\bar\phi_\alpha\phi_\alpha v_{12}
\tag{3.36}
\end{equation}
\begin{equation}
\frac{1}{2e^2}\int \mathrm{d}^2y\,(v_{12}^2+D^2)
= \frac{1}{2e^2}\int \mathrm{d}^2y\,(v_{12}+D)^2
- \frac{1}{e^2}\int \mathrm{d}^2y\,D v_{12}
\tag{3.37}
\end{equation}
以及 $D=-e^2\left(\sum_\alpha Q_\alpha\bar\phi_\alpha\phi_\alpha-r\right)$，可得：若 $D+v_{12}=(D_1+iD_2)\phi_\alpha=0$，则
\begin{equation}
\int \mathrm{d}^2y\left(\sum_\alpha D_i\bar\phi_\alpha D^i\phi_\alpha
+ \frac{1}{2e^2}(v_{12}^2+D^2)\right)
= -r\int \mathrm{d}^2y\, v_{12}.
\tag{3.38}
\end{equation}
% [Tutor Note: 这里是 Bogomolnyi 型“完成平方”：关键对易子是 $[D_1,D_2]\phi_\alpha=-iQ_\alpha\phi_\alpha v_{12}$（符号取决于 §2 中 $v_{mn}$ 与协变导数的约定，务必核对），于是 $(D_1-iD_2)(D_1+iD_2)\phi_\alpha=\left(D_iD^i+Q_\alpha v_{12}\right)\phi_\alpha$，分部积分后即得 (3.36)。(3.37) 只是恒等式 $(v_{12}+D)^2=v_{12}^2+D^2+2Dv_{12}$；把 (3.36) 与 (3.37) 相加并代入 $D$ 的表达式，交叉项恰好消去，左端成为正定泛函——这正是瞬子方程 (3.35) 的 BPS 界。]

上式左端在模去 $\sigma=F_\alpha=0$、并略去 theta 项之后，与该理论拉格朗日量的玻色部分一致。于是式 (3.35) 的解的作用量为
\begin{equation}
L=-2\pi i tN
\tag{3.39}
\end{equation}
其中 $N$ 是瞬子数
\begin{equation}
N=-\frac{1}{2\pi}\int \mathrm{d}^2y\, v_{12},
\tag{3.40}
\end{equation}
而 $t=ir+\theta/2\pi$；把 $r$ 代换为 $-it$ 计入了作用量中的 theta 项。
% [Tutor Note: 校验符号：$t=ir+\theta/2\pi$ 时 $-2\pi itN=-2\pi i(ir+\theta/2\pi)N=2\pi rN-i\theta N$，故权重 $e^{-L}=e^{-2\pi rN+i\theta N}$：瞬子按 $e^{-2\pi rN}$ 被压低、并带上 $e^{i\theta N}$ 的相位；由 $N=-(1/2\pi)\int v_{12}$ 知 theta 项 $(\theta/2\pi)\int v_{12}=-\theta N$，相位一致。]

% ==== 原文 PDF 第 38 页（印刷页 38）====
\label{op:38}
由于 (3.38) 的左端是正定的，可知 (3.35) 的解必然满足：若 $r>0$ 则 $N\ge0$，若 $r<0$ 则 $N\le0$。于是任何拓扑可观测量的求值所用的“瞬子展开”就是按瞬子的展开，其一般形式为
\begin{equation}
\sum_{k\ge0}a_k\exp(2\pi ikt),
\tag{3.41}
\end{equation}
适用于 $r\to\infty$；或者按反瞬子的展开，其一般形式为
\begin{equation}
\sum_{k\ge0}b_k\exp(-2\pi ikt),
\tag{3.42}
\end{equation}
适用于 $r\to-\infty$。

一个标准的消没定理（负次数的线丛不可能有非零全纯截面）指出：若 $(D_1+iD_2)\phi_\alpha=0$，则
\begin{equation}
\phi_\alpha=0\quad\text{unless}\quad \operatorname{sign}(Q_\alpha)=\operatorname{sign}(N).
\tag{3.43}
\end{equation}
% [Tutor Note: 推导补足：$(D_1+iD_2)\phi_\alpha=0$ 即 $\bar\partial_A\phi_\alpha=0$，故 $\phi_\alpha$ 是线丛 $\mathcal{O}(Q_\alpha N)$ 的全纯截面（在 $\Sigma\cong\mathbf{CP}^1$ 上），其次数为 $\deg=Q_\alpha N$；次数为负的线丛没有非零全纯截面，故必须有 $Q_\alpha N\ge0$（或 $\phi_\alpha\equiv0$），这正是 (3.43)。注意这是在复规范变换下把 $v_{12}$ 化为常数后，由 (3.36) 得到的。]

这可由 (3.36) 推出：只要利用 $(D_1+iD_2)\phi=0$ 在复规范变换（见下文）下的不变性，把 $v_{12}$ 取为常数即可。

下面更仔细地考察瞬子与反瞬子的结构，限于与 C-Y/L-G 对应最简单形式相关的情形：有 $n$ 个荷为 $1$ 的手征超场 $S_i$，以及一个荷为 $-n$ 的手征超场 $P$，而超势为 $W=PG(S_i)$。

\subsubsection*{$r\ll0$}

$r\ll0$ 的情形较为简单，故先考虑它。由 (3.43) 知 $s_i=0$，但 $p\ne0$。$s_i$ 为零保证了 $F_\alpha=0$。剩下的方程是
\begin{equation}
\begin{aligned}
(D_1+iD_2)p &= 0\\
v_{12} &= e^2(-n|p|^2-r).
\end{aligned}
\tag{3.44}
\end{equation}
这就是 Nielsen--Olesen 阿贝尔涡旋线的方程，其解的定性性质是众所周知的。
% [背景补充: Nielsen--Olesen 涡旋是阿贝尔希格斯模型中拓扑稳定的管状解：标量场在横截面上绕涡旋一周相位改变 $2\pi$、在中心为零，规范场相应地携带磁通；这里 $p$ 场扮演希格斯场、$v_{12}$ 是其磁通密度，故 (3.44) 描述的就是这种“涡旋线”。]

% ==== 原文 PDF 第 39 页（印刷页 39）====
\label{op:39}
$p$ 场是有质量的，其质量正比于 $e\sqrt r$，康普顿波长为 $\lambda\sim1/(e\sqrt r)$。由于 $p$ 场的荷为 $-n$，式 (3.44) 的解可以具有瞬子数 $-k/n$，其中 $k\ge0$ 为任意整数。在瞬子数为 $-k/n$ 的反瞬子场中，$p$ 场在 $k$ 个（一般互异的）点 $x_1\ldots x_k\in\Sigma$ 处为零，反瞬子场的能量密度集中在距这些点 $\lambda$ 量级的距离以内。对探测远长于 $\lambda$ 的距离的低能观测者而言，该场看起来集中在 $k$ 个点上；它看起来像是 $k$ 个点状反瞬子的叠加，每个的瞬子数为 $-1/n$。
% [原文疑误: 本段讨论 $r\ll0$，此时 $\sqrt r$ 为虚数；按 (3.44) 的 $v_{12}=e^2(-n|p|^2-r)$，$p$ 的质量平方应正比于 $-r$，质量与康普顿波长中的根号内应为 $-r$（即 $e\sqrt{-r}$、$\lambda\sim1/(e\sqrt{-r})$）。此处照印出的 $\sqrt r$ 抄录，未改公式。]

在某一个 $x_a$ 附近，$p$ 场有单零点，它绕该点一周相位改变 $2\pi$。荷为 $1$ 的 $s_i$ 于是相位改变 $-2\pi/n$。因此，从无质量场 $s_i$ 的低能有效理论来看，就好像在 $k$ 个点 $x_a$ 处插入了“扭场”（twist field），$s_i$ 绕扭场一周相位改变 $-2\pi/n$。当然，在整体情形下，在紧黎曼面 $\Sigma$ 上 $s_i$ 的相位总改变必须是 $2\pi$ 的整数倍，故 $k$ 必须被 $n$ 整除。这是意料之中的；它正是在存在荷为 $1$ 的场时瞬子数的标准量子化。
% [Tutor Note: 物理图像——扭场插入正是 Landau--Ginzburg 轨形（orbifold）中 $\mathbb{Z}_n$ 扭扇区的场：$p$ 的零点给出一个绕数，$s_i$ 绕它一周获得 $-2\pi/n$ 的相位，故每一个这样的零点等价于在 $s_i$ 的 CFT 中插入一个分数荷的扭场。$k$ 必须被 $n$ 整除，正是扭场总分数荷必须使 $s_i$ 单值的要求。]

于是可以看出 $r\ll0$ 的 Landau--Ginzburg 理论中参数 $r,\theta$ 的诠释。这里存在扭场的气体，其化学势使得含 $k$ 个扭场的振幅获得因子
\begin{equation}
\exp(-2\pi itk/n).
\tag{3.45}
\end{equation}

对于真空振幅，或插入了普通（未扭）场的振幅，气体中扭场的数目 $k$ 必须被 $n$ 整除。然而，在 Landau--Ginzburg 模型的理论中已知，$A$ 模型在扭扇区中有 BRST 不变的可观测量。含比如说 $s$ 个扭场的关联函数只收到那些使 $k+s$ 被 $n$ 整除的 $k$ 值的贡献。对给定的 $s$，令 $k_0$ 为使 $k_0+s$ 被 $n$ 整除的最小非负整数。
% [跨页接续: 原文第39页末至第40页首]

% ==== 原文 PDF 第 40 页（印刷页 40）====
\label{op:40}
于是含 $s$ 个扭场的关联函数 $\langle T_1\ldots T_s\rangle$ 当 $r\to-\infty$ 时正比于
\begin{equation}
\exp(-2\pi itk_0/n)
\tag{3.46}
\end{equation}
并且特别地，除非 $k_0=0$，它在 $r=-\infty$ 处为零。这意味着在 $r=-\infty$ 处存在一条在其他情形下并不存在的选择定则：除非 $s$ 被 $n$ 整除，$\langle T_1\ldots T_s\rangle=0$。这条选择定则对应于一个仅在 $r=-\infty$ 处成立的量子 $\mathbb{Z}_n$ 对称性。
% [Tutor Note: 推导补足：把 $t=ir+\theta/2\pi$ 代入 (3.46)，得 $\exp(-2\pi itk_0/n)=\exp(2\pi rk_0/n)\cdot\exp(-i\theta k_0/n)$；当 $r\to-\infty$ 时第一个因子以指数速度趋于零，故一切 $k_0\ne0$ 的振幅消失，只有 $k_0=0$（即 $s$ 被 $n$ 整除）的振幅存活。这种“仅在无穷远处涌现”的对称性正是量子 $\mathbb{Z}_n$ 对称性；它与 $\theta$ 角的周期性 $\theta\to\theta+2\pi$ 相容。]

概括地说，要强调：对 $r\ll0$，低能理论是 Landau--Ginzburg 理论，它没有（光滑的）瞬子。光滑瞬子解的构造依赖于有质量场；这些瞬子解在基本层面是光滑的，但在低能理论中看起来像点状奇点。

\subsubsection*{$r\gg0$}

反之若 $r\gg0$，所处理的是瞬子而非反瞬子；(3.43) 迫使 $p=0$。在超势为 $W=PG(S_i)$ 的情形，需要满足的方程是
\begin{equation}
\begin{aligned}
(D_1+iD_2)s_i &= 0\\
G(s_i) &= 0\\
v_{12} &= e^2\Big(\sum_i|s_i|^2-r\Big).
\end{aligned}
\tag{3.47}
\end{equation}

下面将分析这些方程的含义，但首先，这里先不加任何数学形式化地陈述将要得到的结果。$r\gg0$ 的低能理论是靶空间为 Calabi--Yau 超曲面 $X\subset\mathbf{CP}^{n-1}$ 的 $\sigma$ 模型。当然应当期望找到该 $\sigma$ 模型的瞬子，即 $\Sigma$ 到 $X$ 的全纯映射。然而（正如 $r\ll0$ 的情形），还会找到额外的瞬子：它们在基本层面是完全光滑的客体（有质量场存在），但在有效低能理论中看起来像奇异客体。

% ==== 原文 PDF 第 41 页（印刷页 41）====
\label{op:41}
这些额外瞬子的结构在正比于有质量场康普顿波长的尺度上变化（其余结构则变化得更平缓）。
% [跨页接续: 原文第40页末至第41页首]

这些额外“奇异”瞬子的存在意味着：本文所研究的线性 $\sigma$ 模型的瞬子展开，不能指望与它在低能下归结为的共形不变非线性 $\sigma$ 模型的标准瞬子展开相重合。不过，通常的重整化理论告诉我们两者如何相关。应当考虑把点状瞬子结构（连同有质量场的所有其他效应）积掉，以得到无质量场的有效理论。该理论必定是某种 $\sigma$ 模型，带有某些有效值 $r$ 与 $\theta$，以及某个有效的定义方程 $G=0$。（其余参数被认为在重整化群意义下是“无关”的。）由于 $A$ 模型的可观测量不依赖于特定的多项式 $G$，积掉点状瞬子的效应应当只是改变 $r$ 与 $\theta$ 的有效值。因此，线性 $\sigma$ 模型的瞬子展开，应当与非线性 $\sigma$ 模型通常的瞬子展开相差一个变量 $t$ 的重定义。

现在转入式 (3.47) 的数学分析。显然，这些方程在规范变换
\begin{equation}
s_i\to e^{i\lambda}s_i,\qquad v\to v-\mathrm{d}\lambda,
\tag{3.48}
\end{equation}
（$\lambda$ 为实参数）下不变。实际上容易看出，式 (3.47) 的前两个方程即使在 $\lambda$ 为复参数时也在 (3.48) 下不变。一个非常优美的事实 [32] 是：式 (3.47) 的最后一个方程可以解释为固定前两个方程的复规范不变性的条件。换言之，前两个方程的解空间模去复规范变换，与三个方程联立后的解空间模去实规范变换是相同的。（这可以看作下一节所总结的复商与辛商之间关系的一个无穷维类比。）
% [Tutor Note: 这正是 Hitchin--Kobayashi 对应（Donaldson--Uhlenbeck--Yau 定理）的无穷维版本：方程 (3.47) 的前两式是复方程 $\bar\partial_A s=0$ 与 $G(s)=0$（全纯数据），第三式 $v_{12}=e^2(\sum|s_i|^2-r)$ 是矩映射（moment map）条件，它选取复规范轨道中的唯一“规范固定”代表。用几何不变量理论（GIT）的语言：复商 $\mu^{-1}(0)/G_{\mathbb{C}}$ 与辛商 $\mu^{-1}(0)/G$ 在稳定点上一致，下一节将对此作初等说明。]

% ==== 原文 PDF 第 42 页（印刷页 42）====
\label{op:42}
于是可以只研究式 (3.47) 的前两个方程，模去复规范变换。为使分析特别简单，再假设世界面 $\Sigma$ 的亏格为零（这是 $A$ 模型标准应用中的主要情形）。这带来简化，因为 $\Sigma$ 上任何复线丛都是某个 $k$ 的 $\mathcal{O}(k)$。若 $u,v$ 是 $\Sigma\cong\mathbf{CP}^1$ 的齐次坐标，则 $\mathcal{O}(k)$ 的全纯截面不过是 $u,v$ 的、$k$ 次齐次的全纯函数。后文将用同一记号 $\mathcal{O}(k)$ 表示 $\mathbf{CP}^{n-1}$ 上类似定义的线丛。

在复规范变换之下，式 (3.47) 的第一个方程只是说 $s_i$ 是 $\mathcal{O}(k)$ 的全纯截面，其中 $k$ 为瞬子数。于是 $s_i=s_i(u,v)$ 是 $u$ 与 $v$ 的 $k$ 次齐次多项式，具体地
\begin{equation}
s_i(u,v)=s_{i,k}u^k+s_{i,k-1}u^{k-1}v+\ldots+s_{i,0}v^k.
\tag{3.49}
\end{equation}

整体的标度变换
\begin{equation}
s_i(u,v)\to t s_i(u,v),\qquad t\in\mathbb{C}^*
\tag{3.50}
\end{equation}
对应于规范参数为常数的复规范变换；它们是 $\mathcal{O}(k)$ 的同构。式 (3.47) 的第二个方程说
\begin{equation}
G(s_1(u,v),\ldots,s_n(u,v))=0.
\tag{3.51}
\end{equation}

式 (3.47) 的解的模空间，就是满足 (3.51) 的 $k$ 次多项式 $s_1(u,v),\ldots,s_n(u,v)$ 的空间模去 (3.50)。

现在考虑 $\sigma$ 模型的模空间，即 $k$ 次全纯映射 $\Phi:\Sigma\to X$ 的模空间，其中 $X$ 是 $\mathbf{CP}^{n-1}$ 中由 $G=0$ 给出的超曲面。这个 $\sigma$ 模型模空间容易描述。$\Phi$ 的次数为 $k$，按定义意味着 $\Phi^*(\mathcal{O}(1))=\mathcal{O}(k)$。$\mathbf{CP}^{n-1}$ 的齐次坐标 $s_i$ 是 $\mathcal{O}(1)$ 的全纯截面，故 $\Phi^*(s_i)$ 是
% [跨页接续: 原文第42页末至第43页首]

% ==== 原文 PDF 第 43 页（印刷页 43）====
\label{op:43}
$\mathcal{O}(k)$ 的全纯截面。这意味着任何 $k$ 次全纯映射 $\Phi:\mathbf{CP}^1\to\mathbf{CP}^{n-1}$ 都具有形式
\begin{equation}
(u,v)\to(s_1(u,v),\ldots,s_n(u,v))
\tag{3.52}
\end{equation}
其中 $s_1(u,v),\ldots s_n(u,v)$ 是 $k$ 次多项式，且必须满足 (3.51)。由于在 $\mathbf{CP}^{n-1}$ 中不允许 $s_i$ 同时为零，这些多项式必须没有公共零点。反之，若 $s_i$ 满足 (3.51) 且没有公共零点，则公式 (3.52) 描述一个 $k$ 次全纯映射 $\mathbf{CP}^1\to\mathbf{CP}^{n-1}$。

两个模空间之间的关系现在清楚了。每个 $\sigma$ 模型瞬子都来自 (3.47) 的一个解，但 (3.47) 还有额外的解：在其中 $s_i(u,v)$ 对某些不全为零的 $u,v$ 确实有公共零点。这些额外的解在底层的规范场论中是完全光滑的客体，但在 $s_i$ 的公共零点附近，其变化尺度为 $1/(e\sqrt r)$ 量级；这就是引言性评述中所预告的“点状”行为。这些瞬子只在 $s_i$ 的公共零点附近表现出“点状”行为；在其他地方它们看起来就像 $\sigma$ 模型瞬子。

人们也许以为，$s_i$ 有公共零点这一条件只会非常例外地得到满足，从而点状瞬子也许无关紧要。作为看出并非如此的前奏，先把瞬子模空间的维数算出来。式 (3.51) 断言某个 $u,v$ 的 $kn$ 次齐次多项式为零，因此它有展开式 $G(u,v)=a_{kn}u^{kn}+a_{kn-1}u^{kn-1}v+\ldots+a_0v^{kn}$。这 $kn+1$ 个 $a$ 可以显式地写成 $(k+1)n$ 个系数 $s_{i,r}$ 的 $n$ 次多项式。再计入标度关系 (3.50)，若一切都处于一般位置，瞬子模空间的维数预期为 $(k+1)n-(kn+1)-1=n-2$。已知对一般的 $G$，这确实是 $\sigma$ 模型模空间的实际维数。（在靶空间为 $\mathbf{CP}^4$ 中五次三维流形这一通常情形，$n-2=3$，而在除以 $\Sigma$ 上 $SL(2,\mathbb{C})$ 的作用之后维数降为零。）
% [Tutor Note: 这是 Brill--Noether 型的“方程数减未知数”计数：$(k+1)n$ 个系数（每个 $s_i$ 有 $k+1$ 个系数）减去 $kn+1$ 个独立方程、再减去标度（1 个），得 $n-2$。这个维数正是量子上同调/Gromov--Witten 不变量所计数的全纯映射模空间的维数（在除以世界面自同构之前）；除以 $SL(2,\mathbb{C})$ 是去掉世界面 $\mathbf{CP}^1$ 的自同构（三个点），故得到 $n-2-3$ 维，对五次三维流形恰为 0 维，与 $g=0$ 三点函数的计数相符。]

若现在 $s_i(u,v)$ 是一组 $k$ 次多项式，满足 (3.51) 且没有公共零点，则 $\tilde s_i(u,v)=(\alpha u+\beta v)s_i(u,v)$ 是一组 $k+1$ 次
% [跨页接续: 原文第43页末至第44页首；p43 末句 “...are a collection of degree” 在 p44 首行续为 “k+1 polynomials, obeying (3.51) ...”；p44 仅前 9 行属本分片，§4 标题起归 s04-geometry]
% ==== 原文 PDF 第 44 页（印刷页 44；仅前 9 行属本分片）====
\label{op:44}
多项式，满足 (3.51)，并且在 $\alpha u+\beta v=0$ 处有公共零点。在计及 $\alpha,\beta$ 的一个无关标度之后，$\tilde s_i$ 比 $s_i$ 多依赖一个参数；因此，$\sigma$ 模型瞬子依赖 $n-2$ 个参数，而具有一个公共零点的点状瞬子依赖 $n-1$ 个参数。（而具有 $w$ 个公共零点的点状瞬子依赖 $n-2+w$ 个参数。）点状瞬子绝非例外，在把有质量场积掉以得到无质量场的有效作用量时必须考虑它们。正如前面所述，通常的重整化经验提示，它们的效应是诱导变量 $t$ 的一个重参数化。
